%%%%%%%%%%%%%%%%%%%%%%%%%%%%%%%%%%%%%%%%%
% Short Sectioned Assignment LaTeX Template Version 1.0 (5/5/12)
% This template has been downloaded from: http://www.LaTeXTemplates.com
% Original author:  Frits Wenneker (http://www.howtotex.com)
% License: CC BY-NC-SA 3.0 (http://creativecommons.org/licenses/by-nc-sa/3.0/)
%%%%%%%%%%%%%%%%%%%%%%%%%%%%%%%%%%%%%%%%%

% \documentclass[paper=a4, fontsize=11pt]{scrartcl} % A4 paper and 11pt font size
\documentclass[11pt, a4paper]{book}
\usepackage[T1]{fontenc} % Use 8-bit encoding that has 256 glyphs
\usepackage[utf8]{inputenc}
\usepackage{listings} % para insertar código con formato similar al editor
\usepackage[spanish, es-tabla]{babel} % Selecciona el español para palabras introducidas automáticamente, p.ej. "septiembre" en la fecha y especifica que se use la palabra Tabla en vez de Cuadro
\usepackage{url} % ,href} %para incluir URLs e hipervínculos dentro del texto (aunque hay que instalar href)
\usepackage{graphics,graphicx, float} %para incluir imágenes y colocarlas
\usepackage[gen]{eurosym} %para incluir el símbolo del euro
\usepackage{cite} %para incluir citas del archivo <nombre>.bib
\usepackage{enumerate}
\usepackage{hyperref}
\usepackage{graphicx}
\usepackage{tabularx}
\usepackage{booktabs}

\newcommand{\titulo}{MaquiTracker}
\newcommand{\subtitulo}{Servicio de localización, estadística y mantenimiento preventivo para maquinaria agrícola.}
\newcommand{\tituloingles}{MaquiTracker}
\newcommand{\subtituloingles}{Location, statistics and preventive maintenance service for agricultural machinery.}
\newcommand{\autor}{Daniel Pedrosa Montes}
\newcommand{\tutor}{Juan Julián Merelo Guervos}
\newcommand{\departamento}{Departamento de Ingeniería de Computadores, Automática y Robótica}

\usepackage[table,xcdraw]{xcolor}
\hypersetup{
	colorlinks=true,	% false: boxed links; true: colored links
	linkcolor=black,	% color of internal links
	urlcolor=cyan		% color of external links
}

% Configuración tipografía
\renewcommand{\familydefault}{\sfdefault}
\usepackage[scaled=1]{helvet}
\usepackage[helvet]{sfmath}
\everymath={\sf}

% Header y footer
\usepackage{fancyhdr} % Custom headers and footers
\pagestyle{fancyplain} % Makes all pages in the document conform to the custom headers and footers
\fancyhead[L]{} % Empty left header
\fancyhead[C]{} % Empty center header
\fancyhead[R]{\autor} % My name
\fancyfoot[L]{} % Empty left footer
\fancyfoot[C]{} % Empty center footer
\fancyfoot[R]{\thepage} % Page numbering for right footer
%\renewcommand{\headrulewidth}{0pt} % Remove header underlines
\renewcommand{\footrulewidth}{0pt} % Remove footer underlines
\setlength{\headheight}{13.6pt} % Customize the height of the header

\usepackage{titlesec, blindtext, color}
\definecolor{gray75}{gray}{0.75}
\newcommand{\hsp}{\hspace{20pt}}
\titleformat{\chapter}[hang]{\Huge\bfseries}{\thechapter\hsp\textcolor{gray75}{|}\hsp}{0pt}{\Huge\bfseries}
\setcounter{secnumdepth}{4}
\usepackage[Lenny]{fncychap}


\begin{document}

	% Plantilla portada UGR
	\begin{titlepage}
\newlength{\centeroffset}
\setlength{\centeroffset}{-0.5\oddsidemargin}
\addtolength{\centeroffset}{0.5\evensidemargin}
\thispagestyle{empty}

\noindent\hspace*{\centeroffset}\begin{minipage}{\textwidth}

\centering
\includegraphics[width=0.9\textwidth]{logos/logo_ugr.jpg}\\[1.4cm]

\textsc{ \Large TRABAJO FIN DE GRADO\\[0.2cm]}
\textsc{ GRADO EN INGENIERIA INFORMATICA}\\[1cm]

{\Huge\bfseries \titulo \\}
\noindent\rule[-1ex]{\textwidth}{3pt}\\[3.5ex]
{\large\bfseries \subtitulo }
\end{minipage}

\vspace{2.5cm}
\noindent\hspace*{\centeroffset}
\begin{minipage}{\textwidth}
\centering

\textbf{Autor}\\ \autor \\[2.5ex]
\textbf{Director}\\ {\tutor}\\[2cm]
\includegraphics[width=0.3\textwidth]{logos/etsiit_logo.png}\\[0.1cm]
\textsc{Escuela Técnica Superior de Ingenierías Informática y de Telecomunicación}\\
\textsc{---}\\
Granada, Junio de 2024
\end{minipage}
\end{titlepage}


	% Plantilla prefacio UGR
	\thispagestyle{empty}

\begin{center}
{\large\bfseries \titulo \\ \subtitulo }\\
\end{center}
\begin{center}
\titulo \\
\end{center}

%\vspace{0.7cm}

\vspace{0.5cm}
\noindent\textbf{Palabras clave}: \textit{software libre}
\vspace{0.7cm}

\noindent\textbf{Resumen}\\
	

\cleardoublepage

\begin{center}
	{\large\bfseries \tituloingles \\ \subtituloingles}\\
\end{center}
\begin{center}
	\autor\\
\end{center}
\vspace{0.5cm}
\noindent\textbf{Keywords}: \textit{open source}, \textit{floss}
\vspace{0.7cm}

\noindent\textbf{Abstract}\\


\cleardoublepage

\thispagestyle{empty}

\noindent\rule[-1ex]{\textwidth}{2pt}\\[4.5ex]

D. \textbf{\tutor}, Profesor(a) del \departamento.

\vspace{0.5cm}

\textbf{Informo:}

\vspace{0.5cm}

Que el presente trabajo, titulado \textit{\textbf{\titulo}},
ha sido realizado bajo mi supervisión por \textbf{\autor}, y autorizo la defensa de dicho trabajo ante el tribunal
que corresponda.

\vspace{0.5cm}

Y para que conste, expiden y firman el presente informe en Granada a Junio de 2018.

\vspace{1cm}

\textbf{El/la director(a)/es: }

\vspace{5cm}

\noindent \textbf{\tutor}

\chapter*{Agradecimientos}






	% Índice de contenidos
	\newpage
	\tableofcontents

	% Índice de imágenes y tablas
	\newpage
	\listoffigures

	% Si hay suficientes se incluirá dicho índice
	\listoftables 
	\newpage

	% Introducción 
	\chapter{Introducción}

Este Trabajo de Fin de Grado consiste en construir una solución software
capaz de emitir, monitorizar y registrar una serie de eventos en un
sistema de control hidráulico embebido. 

Este proyecto es software libre, y está liberado con
la licencia \textbf{GNU Affero General Public License} \cite{agplv3}.

\section{Motivación}

En España el mercado de los cultivos leñosos supone un ??? del PIB
anual, siendo un sector que emplea a más de ??? personas. [cita] Entre estos
predominan el cultivo de olivos, almendros, pistachos, cítricos y
viñedos. [cita] 

La explotación de estas plantaciones se realiza mayoritariamente por personas
físicas y no por empresas. [cita] La profesionalización de los propietarios
es creciente, pero baja. [cita] A su vez ocurre que no existe una renovación
de los jefes de explotación, con un porcentaje de gestores menores de 45 años
inferior al 40\%. La formación reglada específica de los jefes en el caso de
las explotaciones agrarias es inferior al 25\% en todas las comarcas españolas,
llegando a menos del ??? en aquellas regiones con mayor
explotación de este tipo de cultivos. [cita] Encontramos la mayor superficie
de cultivos leñosos en la mitad sur de la península, destacando Andalucía,
Castilla la Mancha y Extremadura. [cita INE https://storymaps.arcgis.com/stories/6fa8c26ecbde4fd99a0446e874c64898]

Desgraciadamente, el bajo nivel de automatización y de inversión por parte
de los propietarios de las explotaciones agrarias hace que la rentabilidad
de estos haya disminuido en los últimos años. [cita]
El sector de la
mecanización de este tipo de cultivos ha crecido en los últimos años de una
forma lenta y con algunas pausas, dado el alto coste de la maquinaria, la
falta de emprendimiento en el sector y el hecho de que la mayoría de las
explotaciones son de pequeño tamaño. [cita] 

Dos terceras partes del
coste de este tipo de cultivos van destinados a la recolección y a la poda.
[cita https://centroliva.com/mecanizacion-de-la-recoleccion-del-olivar-desde-la-investigacion-hasta-la-industria/]
Entre las soluciones mecanizadas
para la recolección encontramos principalmente los vibradores y los paraguas,
mientras que para la poda encontramos las tijeras eléctricas, las sierras y
las trituradoras. [cita]

[ Crear o extraer gráficos del INE que ayuden a entender los datos aquí expuestos ]

Son muchos los jefes de explotación que deciden externalizar la recolección
de sus cultivos a empresas especializadas. [cita?] Ocurre un problema en este
caso, y es que el jefe de explotación deja de tener la certeza de que la
recolección se ha realizado de forma correcta, y de que la cantidad de producto
recolectado es la que se espera. Este problema sucede también en otros
contextos, como el de grandes fincas profesionalizadas.


\section{Objetivo}

Este proyecto busca desarrollar una solución software que permita
registrar, exportar y visualizar información sobre la recolección de cultivos
leñosos con paraguas vibradores según el criterio de nuestro cliente,
con la restricción de que el registro y el exportado pueda ejecutarse en
una plataforma hardware ya distribuida en el mercado.

Entre otros objetivos secundarios se encuentran:

\begin{enumerate}
    \item Automatización de los procesos de desarrollo y despliegue de la solución, como es
    la generación de documentación, la ejecución de pruebas, la generación de librerías
    y el despliegue en la plataforma hardware.
    \item Escribir código que cumpla los estándares de la industria para maquinaria pesada,
    asegurando que no puedan ocurrir situaciones inesperadas en tiempo de ejecución que puedan
    comprometer la integridad del código de control, que se ejecutará sobre la misma plataforma
    hardware.
    \item Generar una documentación adecuada para que cualquier desarrollador pueda entender,
    mantener y ampliar este proyecto en el futuro.
\end{enumerate}

\section{Estructura del documento}

[ Explica qué secciones tiene el documento y qué desarrollan ]

[ Explica el por qué de la elección ]


	% Descripción del problema y hasta donde se llega
	\chapter{Descripción del problema}



	% Estado del arte
	% 	1. Crítica al estado del arte
	% 	2. Propuesta
	\chapter{Estado del arte}

El software libre y sus licencias \cite{agplv3} ha permitido llevar a cabo una expansión del
aprendizaje de la informática sin precedentes.

	
	\chapter{Planificación}

\section{Metodología utilizada}

[ Describe Trunk Based Development ]

[ Describe Continuous Delivery ]

[ Explica el por qué de la elección ]

[ ¿Anexo sobre la historia de las distintas metodologías y
filosofías de desarrollo de software? ]

Para llevar a cabo este producto software intentaremos seguir algunas
pautas que buscarán acortar nuestros tiempos de \textit{feedback}.

\subsection{¿Por qué acortar los tiempos de \textit{feedback}?}

Básicamente porque \textbf{cuanto antes detectamos que nos estamos equivocando
antes podremos corregirlo}, y antes podremos seguir adelante.

Pensemos en la situación de un enamorado sin experiencia que durante meses dedica su
tiempo, su energía y su creatividad a componerle una canción a una chica.
Durante este tiempo, este individuo mejora de forma objetiva la
melodía, letra y sensaciones que la canción genera en sí mismo
y en sus amistades. Ahora es un mejor compositor, tiene un mejor manejo
de los instrumentos que utiliza e incluso ha mejorado su
relación con sus amigos y su familia.

El chico piensa que así puede despertar o avivar un sentimiento recíproco
en la chica. Sin embargo, cuando se presenta la canción,
la realidad se impone. Entre otros muchos factores,
la falta de sincronización entre los sentimientos de ella y de él causan
de forma instintiva un rechazo en ella.

Esta analogía \textbf{la mostramos para ejemplificar lo absurdas e idealistas
que son estas situaciones que se dan en otros campos de nuestra vida.}
Productos excelentes desde un punto de vista técnico que se
congelan porque no se discutió un precio objetivo con
el cliente en una fase temprana del desarrollo, ideas felices
por parte de \textit{stakeholders} sobre qué debe de
ser capaz de hacer el software que terminan en años de desarrollo
y con objetivos sin cumplir, etc.

Si el chico le hubiese propuesto de quedar de una forma casual
habría llegado a conocerla mejor\footnote{¡O no! A lo mejor ella no
quiera quedar con él. Al chico le interesa
conocer esta realidad lo antes posible.}. Podría haberse dado cuenta que
ella no está acostumbrada a salir mucho de casa, que buscando sitios
interesantes con los que pasar tiempo con ella tal vez podría ganar su
confianza, y su cariño. Podría haber probado esta hipótesis, y cambiar
su estrategia según el \textit{feedback} que recibiría. Podría en este tiempo que
pasan juntos darse cuenta que no es la chica adecuada para él. Podría
incluso haberse dado cuenta que, la idea de la canción, podría ser buena
bajo unas determinadas circunstancias.

Los autores de \cite{accelerate} divulgan cómo afectan los tiempos de feedback
a los productos software. Nos centramos en los capítulos dedicados al impacto de
la entrega continua y al impacto de la entrega continua%
\footnote{%
Superconjunto de la integración continua. Básicamente consiste en la automatización
del proceso de desarrollo. \cite{ModernSoftwareEngineering} y \cite{minimumViableCD}
son buenos recursos para aprender más acerca de estos conceptos.%
}%
en la calidad. Según escriben, \textbf{existe una correlación estadística alta entre aquellos
equipos capaces de desplegar a producción a demanda y 
con mecanismos que ofrecen de forma clara un \textit{feedback} rápido en
la calidad y en la capacidad del software de desplegarse y entre el rendimiento del
equipo.}

El equipo DORA\footnote{DevOps Research and Assessment.} coincide en estos resultados.
En su publicación \cite{EliteDevOps} explican cuatro métricas que pueden medir el
rendimiento de una organización. Según explican en la entrada de blog citada,
la frecuencia de despliegue, el tiempo que tarda un commit desde que se sube
al servidor hasta que se lleva a producción,
el porcentaje de despliegues que causan un error en producción y el tiempo
que tarda una organización en recuperarse de un error son clave \textbf{para medir la
velocidad y la estabilidad en el desarrollo}.

\subsection{Productos para disminuir el tiempo de realimentación}

\begin{enumerate}
    \item \textbf{Automatización de los procesos de desarrollo y despliegue de la solución}, como es
    la generación de documentación, la ejecución de pruebas, la generación de librerías
    y el despliegue en la plataforma hardware, \textbf{de forma que minimicemos el tiempo de \textit{feedback}}
    una vez realicemos un cambio y \textbf{de forma que tengamos una única fuente de verdad} acerca de cómo
    validar y cómo desplegar.
\end{enumerate}

\subsection{Productos para mejorar la calidad del producto}

\begin{enumerate}
    \item \textbf{Pruebas software unitarias para cada funcionalidad implementada} que aseguren que no
    puedan ocurrir situaciones inesperadas en tiempo
    de ejecución que puedan comprometer la integridad del código de control, que se ejecutará sobre
    la misma plataforma hardware.
    \item \textbf{Documentación adecuada para que cualquier desarrollador pueda entender,
    mantener y ampliar este proyecto} en el futuro, de forma que maximicemos el atractivo
    del proyecto y de forma que minimicemos la frustración de un contribuyente potencial
    al proyecto de software libre.
\end{enumerate}

\section{Temporización}

[ Describe una planificación del tiempo de desarrollo del proyecto ]

[ Explica el por qué de la elección ]

\section{Seguimiento del desarrollo}

[ Describe la plantilla de costes no recurrentes asociada al proyecto ]

[ Explica el por qué de la elección ]


	% Análisis del problema
	% 1. Análisis de requisitos
	% 2. Análisis de las soluciones
	% 3. Solucion propuesta
	% 4. Análisis de seguridad
	\chapter{Análisis del problema}
 


	% Desarrollo bajo sprints: 
	% 	1. Permitir registros y login de usuarios
	% 	2. Desarrollo del sistema de incidencias
	% 	3. Desarrollo del sistema de denuncias administrativas y accidentes
	% 	4. Desarrollo del sistema de croquis
	%   5. Instalación de la aplicación de manera automática
	\chapter{Implementación}

\section{Estimación de costes}

\section{Almacenamiento de histórico}

\section{Recuperación de histórico}


	% Presupuesto

	% Conclusiones
	\chapter{Conclusiones y trabajos futuros}



	% Trabajos futuros


	
	\newpage
	\bibliography{bibliografia}
	\bibliographystyle{plain}
	
\end{document}

